\subsubsection{Speciation and extinction dynamics in New World monkeys}

\begin{figure}[htpb]
    \includegraphics[width=1\textwidth]{figures/platyrrhine/summary.pdf}
    \caption{Diversification analysis of 100 New World monkeys (platyrrhines) phylogenies using the BELLA framework. The analysis estimates time-varying speciation rates (\( \lambda \), top row), extinction rates (\( \mu \), middle row), and net diversification rates (\( d = \lambda - \mu \), bottom row) as functions of time and a discretized body-mass trait evolving along the phylogeny. (a) The estimated evolution of the discretized body-mass trait across the maximum clade credibility tree. (b) The median posterior estimates across time and body-mass categories for each macroevolutionary rate (full distributions are presented in Fig. \ref{fig:platyrrhine:ribbons}). (c) SHAP feature importance distributions for each predictor and macroevolutionary rate.}%
    \label{fig:platyrrhine}
\end{figure}

We applied our BELLA framework to explore macroevolutionary dynamics in the New World monkeys, to estimate their speciation and extinction rates from a set of 100 phylogenetic trees of 87 extant and 34 extinct species from a previous Bayesian analysis~\citep{silvestro2019early}.
We used time (binned into geological stages) and body mass (discretized into four categories) as predictors of rate variation (Fig.~\ref{fig:platyrrhine}a).

Overall, the model captures contrasting effects of time and body mass on the evolutionary dynamics (Fig.~\ref{fig:platyrrhine}b).
The speciation rate ($\lambda$) remains largely constant throughout the entire process and across all body-mass categories, showing little evidence for substantial temporal or trait-dependent variation.
Minor fluctuations are present, but no pronounced shifts comparable to those observed in extinction or diversification rates emerge among the different body-size groups.

Extinction rates ($\mu$) exhibit stronger temporal variation and clearer interactions with body mass.
All body-mass categories experienced an increase in extinction rates around 10 Ma, although the magnitude of this increase differed among groups.
Smaller-bodied species (body-mass category 0) show an increase of approximately 80\%, whereas larger-bodied species (body-mass categories 1, 2 and 3) display a stronger rise of roughly 150\%.
However, posterior median estimates indicate substantial uncertainty around these extinction peaks, making the intensity of the event less clearly supported (Fig.~\ref{fig:platyrrhine:ribbons}).
Following this interval, extinction rates decline across all body-mass categories from the late Miocene through the Recent.

The diversification rates ($d = \lambda - \mu$) reflect these temporal dynamics and their interaction with body mass.
All body-mass groups exhibit a decline in diversification around 10 Ma followed by a subsequent recovery toward the present, producing an overall inverse pattern relative to extinction rates.
In contrast to the extinction estimates, uncertainty around diversification trajectories is considerably more constrained, providing stronger support for a transient decrease followed by a progressive increase in diversification rates through time (Fig.~\ref{fig:platyrrhine:ribbons}).
This decline is particularly pronounced in larger-bodied species, whereas smaller-bodied taxa maintain comparatively more stable diversification dynamics.

The SHAP analysis of feature importance revealed that time consistently had greater importance than body mass across all rates, with the strongest dominance observed for extinction and diversification dynamics (Fig.~\ref{fig:platyrrhine}c). Nevertheless, body mass still contributes to explaining differences among lineages, particularly by separating the stronger extinction and diversification shifts observed in larger-bodied species from the comparatively weaker responses detected in the smallest taxa.